\documentclass[a4paper,11pt]{article}
\usepackage[a4paper,margin=2cm]{geometry}
\usepackage[T1]{fontenc}
\usepackage[utf8]{inputenc}
\usepackage[ngerman,english]{babel}
\usepackage{lmodern}
\usepackage{microtype}
\usepackage{xcolor}
\usepackage{parskip}
\usepackage{enumitem}
\usepackage[most]{tcolorbox}
\usepackage{titlesec}
\usepackage[hidelinks]{hyperref}
\definecolor{HeaderBlueDark}{RGB}{70,110,170}
\definecolor{RowBlueLight}{RGB}{235,243,255}
\definecolor{RowBlueDark}{RGB}{220,233,250}
\definecolor{SoftGray}{RGB}{90,90,90}
\setlength{\parindent}{0pt}
\setlist[itemize]{leftmargin=1.4em,itemsep=0.35em,topsep=0.35em}
\linespread{1.06}
\titleformat{\section}[block]{\Large\bfseries\color{HeaderBlueDark}}{}{0pt}{}
\titlespacing{\section}{0pt}{1.3em}{0.8em}
\titleformat{\subsection}[block]{\large\bfseries\color{HeaderBlueDark}}{}{0pt}{}
\titlespacing{\subsection}{0pt}{1.0em}{0.6em}
\newtcolorbox{exbox}{enhanced,breakable,colback=RowBlueLight,colframe=RowBlueDark,boxrule=0.4pt,arc=1.5mm,left=1.2mm,right=1.2mm,top=1mm,bottom=1mm,title={Beispiele \textnormal{(Examples)}},fonttitle=\bfseries,colbacktitle=RowBlueDark,coltitle=black}
\NewDocumentEnvironment{grammarentry}{mm+b}{%
\begin{tcolorbox}[enhanced,breakable,colback=white,colframe=HeaderBlueDark,boxrule=0.6pt,arc=1.5mm,left=1.5mm,right=1.5mm,top=1mm,bottom=1mm,colbacktitle=HeaderBlueDark,coltitle=white,fonttitle=\bfseries,title={#1\hfill\normalfont\footnotesize #2}]
#3
\end{tcolorbox}}{}
\newcommand{\GermanText}[1]{\textbf{Deutsch: }#1\par}
\newcommand{\EnglishText}[1]{{\small\color{SoftGray}\textbf{English: }#1\par}}
\newcommand{\ExampleItem}[2]{\item \textbf{#1}\\{\small\color{SoftGray}#2}}


\title{Relativwort: \glqq wovon\grqq{}}

\date{}

\begin{document}

\maketitle

\tableofcontents

\newpage

\section{Grundbedeutung und Aufbau}

\begin{grammarentry}{wovon}{Pronominaladverb / Relativadverb}

\GermanText{\textbf{wovon} besteht aus \textbf{wo(r)-} und der Präposition \textbf{von}. Es kann in einem Relativsatz auf eine Sache, einen unbestimmten Inhalt oder einen Sachverhalt verweisen.}

\EnglishText{\textbf{wovon} consists of \textbf{wo(r)-} and the preposition \textbf{von}. In a relative clause, it can refer to a thing, indefinite content, or a situation.}

\GermanText{Ungefähre englische Bedeutung: \glqq of/from which; what ... of/from\grqq{}.}

\EnglishText{Approximate English meaning: \glqq of/from which; what ... of/from\grqq{}.}

\end{grammarentry}

\section{Keine Kasusendung}

\begin{grammarentry}{Pronominaladverb}{nicht deklinierbar}

\GermanText{\textbf{wovon} selbst hat keinen Nominativ, Akkusativ, Dativ oder Genitiv. Die Präposition bleibt aber Teil der Bedeutung und gehört oft fest zu einem Verb oder Ausdruck.}

\EnglishText{\textbf{wovon} itself has no nominative, accusative, dative, or genitive. The preposition remains part of the meaning and is often fixed to a verb or expression.}

\GermanText{Typische Verbindungen: \textbf{träumen von}; \textbf{abhängen von}.}

\EnglishText{Typical combinations: \textbf{träumen von}; \textbf{abhängen von}.}

\end{grammarentry}

\section{Beispiele}

\begin{exbox}

\begin{itemize}

\ExampleItem{Das ist etwas, wovon ich lange geträumt habe.}{That is something I have dreamed of for a long time.}

\ExampleItem{Es gibt vieles, wovon der Erfolg abhängen kann.}{There are many things on which success can depend.}

\end{itemize}

\end{exbox}

\section{Bezug auf Personen}

\begin{grammarentry}{Wichtige Grenze}{bei Personen andere Form}

\GermanText{Bei Personen benutzt man normalerweise nicht \textbf{wovon}, sondern \textbf{von + Relativpronomen}, zum Beispiel \textbf{mit dem / mit der} oder \textbf{für den / für die}.}

\EnglishText{With people, German normally does not use \textbf{wovon}. Instead, it uses \textbf{von + relative pronoun}, for example \textbf{mit dem / mit der} or \textbf{für den / für die}.}

\end{grammarentry}

\section{Konkretes Nomen oder unbestimmter Bezug}

\begin{grammarentry}{Stilregel}{für B1 hilfreich}

\GermanText{Nach \textbf{alles, etwas, nichts, das} sind wo(r)-Formen besonders typisch. Bei einem konkreten Nomen ist \textbf{Präposition + der/die/das} häufig die klarere Standardsprache.}

\EnglishText{After \textbf{alles, etwas, nichts, das}, wo(r)-forms are especially typical. With a concrete noun, \textbf{preposition + der/die/das} is often the clearer standard-language form.}

\GermanText{Muster: \textbf{etwas, wovon ...}; aber oft \textbf{das konkrete Nomen, von dem/der ...}.}

\EnglishText{Pattern: \textbf{etwas, wovon ...}; but often \textbf{the concrete noun, von dem/der ...}.}

\end{grammarentry}

\section{Wortstellung und Komma}

\begin{grammarentry}{Nebensatzregel}{Verb am Ende}

\GermanText{Der Relativsatz ist ein Nebensatz. Das konjugierte Verb steht normalerweise am Ende und der Satz wird durch Kommas abgetrennt.}

\EnglishText{The relative clause is a subordinate clause. The conjugated verb normally stands at the end and the clause is separated by commas.}

\end{grammarentry}

\section{Merksatz}

\begin{grammarentry}{wovon}{kurz und wichtig}

\GermanText{\textbf{wovon} kann relativ verwendet werden, ist aber grammatisch ein \textbf{Pronominaladverb/Relativadverb}, kein deklinierbares Relativpronomen.}

\EnglishText{\textbf{wovon} can be used relatively, but grammatically it is a \textbf{pronominal/relative adverb}, not a declinable relative pronoun.}

\end{grammarentry}

\end{document}